% p.66 — الباب التاسع والعشرون في الاصطلاحات والتركيبات (تتمة: التركيبات)
\mspage{66}
% Layout: four columns = two pairs (label | text), right pair first; continues the
% تركيب examples (the first three rows complete example (19) of p.65).  Circled
% example numbers written as (n).  Printed page number ۶۶ at top left (an older
% number ۵۲ is faintly visible at top right).  Layout: the two pairs are set side by side
% (right pair first) in one `tables` row, as in the lithograph.
\begin{tables}
\begin{sarf}{|r|r|}
فعل & سرق فعل\\
فاعل & زيد فاعل\\
مفعول & القلمَ مفعول\\
 & (20) سُرِقَ الْقَلَمُ\\
فعل مجهول & سرق فعل مجهول\\
نائب فاعل & القلم نائب فاعل\\
 & (21) زَيْدٌ ضَرَبَ عَمْرًا\\
فعل ضمير فاعل & ضَرَبَ فعل ضمير هو فاعل\\
 & (22) ضَرَبْتُ زَيْدًا\\
فعل ضمير فاعل & ضَرَبْتُ فعل ضمير انا فاعل\\
 & (23) اِضْرِبْ زَيْدًا\\
فعل ضمير فاعل & اِضْرِبْ فعل ضمير انتَ فاعل\\
 & (24) أَلْقَلَمُ سُرِقَ\\
فعل ضمير نائب فاعل & سُرِقَ فعل ضمير هو نائب فاعل\\
 & (25) يَا زَيْدُ يَا عَبْدَ اللّٰهِ\\
حرف نداء & يا حرف نداء\\
منادى & زيد منادى عبد الله منادى\\
 & (26) ضَرَبَ زَيْدٌ عَمْرًا مَشْدُوْدًا\\
حال & مشدودا حال\\
ذو الحال & عمرًا ذو الحال\\
 & (27) ذهب زيد الى المدراس وهو راكب\\
\end{sarf}
\begin{sarf}{|r|r|}
جمله حال & وهو راكب جملة حال\\
حالية & واو حالية\\
ذو الحال & زيد ذو الحال\\
 & (28) لَسْتُ أَرْضٰى قَفَصًا وَ اِنْ يَكُنْ مِنْ ذَهَبٍ\\
وصلية & واو وصلية\\
 & (29) مَا جَاءَنِي النَّاسُ اِلَّا زَيْدًا\\
استثناء & الا حرف استثناء\\
مستثنى به & زَيْدًا مُسْتَثْنٰى بِهِ\\
مستثنى منه & الناس مستثنى منه\\
 & (30) رَأَيْتُ أَحَدَ عَشَرَ كَوْكَبًا\\
مُمَيَّز & احد عشر مُمَيَّز\\
تمييز & كوكبا تمييز\\
 & (31) هَيْهَاتَ زَيْدٌ\\
اسم فعل & هيهات اسم فعل\\
فاعل & زَيْدٌ فَاعِلٌ\\
 & (32) هَا زَيْدًا\\
اسم فعل & هَا اسم فعل\\
مفعول & زَيْدًا مفعول\\
 & (33) سَجَدَ الْمَلَائِكَةُ كُلُّهُمْ أَجْمَعُوْنَ\\
تأكيد & كُلُّهُمْ اجمعون : تأكيد\\
مُؤَكَّد & الملائكة مؤكد\\
\end{sarf}
\end{tables}
